\begin{table}[tbp]
\caption{Measurements, models, and the claims they can support.}
\label{tab:claims}
\begingroup\fontsize{9.2}{11.8}\selectfont
\renewcommand{\arraystretch}{1.12}
\begin{tabularx}{\textwidth}{@{}>{\raggedright\arraybackslash}p{25mm}>{\raggedright\arraybackslash}X>{\raggedright\arraybackslash}X@{}}
\toprule
\textbf{Claim} & \textbf{Example in open collaboration} & \textbf{Required connection}\\
\midrule
Description & Recorded participation increased. & Population, coverage, identity, eligible events, and period.\\[5pt]
Measurement & Activity represents a specified account of participation. & Meaning of behaviours, attribution, units, transformation, and aggregation.\\[5pt]
Relational influence & An actor receives a high OpenRank score. & Graph semantics, priors, propagation assumptions, and comparison population.\\[5pt]
Prediction & A procedure anticipates a later outcome. & Test observations unavailable during fitting and selection; a stated target.\\[5pt]
Causal attribution & A recognition arrangement changes collaboration. & Actual exposure, a defensible reference, and identification assumptions.\\[5pt]
Value judgment & The arrangement warrants continued support. & Consequences, burdens, stakeholder requirements, and alternatives.\\
\bottomrule
\end{tabularx}
\endgroup
\tablenote{Book synthesis. Examples state kinds of claim rather than new findings. These claims impose different evidential obligations; they are not an automatic progression from one to the next. Chapters 3--6 develop the corresponding designs, methods, and evidence.}
\end{table}
